\chapter*{\texthebrew{תקציר}}
\addcontentsline{toc}{chapter}{Abstract (Hebrew)}
\markboth{Abstract (Hebrew)}{Abstract (Hebrew)}
\label{sec:abstract-hebrew}
\begin{otherlanguage}{hebrew}
\begin{justify}
עבודה זו משווה עיבוד קלאסי ועיבוד נלמד במערכת תקשורת דו־דילוגית עם ממסר יחיד וללא נתיב ישיר בין המקור ליעד. היא בוחנת מתי ממסר נלמד קטן משפר את ביצועי קווי הבסיס הקלאסיים שמומשו, וכיצד ההשוואה תלויה במידע על הערוץ ובעלות החישוב.

הניסוי הקנוני משתמש באפנון \textenglish{Gray-QPSK} ללא קוד ערוץ, בקישור \textenglish{SISO} עם דעיכת ריילי בלתי־תלויה, ובידיעת ערוץ בצד המקלט. שמונה אסטרטגיות ממסר מושוות בסימולציית מונטה קרלו. ממסרים נלמדים משפרים את הביצועים ביחס להגבר־והעבר (\textenglish{AF}) ביחס אות־רעש נמוך, אך פענח־והעבר סימבולי (\textenglish{DF}) משתווה לממסרים הנלמדים שנבדקו או עולה עליהם החל מ־\textenglish{6 dB}. סריקות גודל משוכפלות מזהות תצורות קטנות העומדות בקריטריוני הירידה בביצועים שנקבעו; הן אינן מוכיחות מנגנון של התאמת־יתר או ארכיטקטורה מזערית אוניברסלית.

ניסויים נפרדים מוסיפים הפרעה בין־סימבולית, אי־ליניאריות של מגבר ואי־ודאות במידע על הערוץ. בערוץ \textenglish{BPSK} הקבוע שנבחר, בעל שלושה מקדמים, ממסר \textenglish{MLP} ממוסגר־חלון בן 170 פרמטרים מפחית במידה ניכרת את שיעור שגיאות הסיביות ביחס ל־\textenglish{AF} ול־\textenglish{DF} בעל סף אפס. גלאי \textenglish{Viterbi MLSE} מותאם־מודל עדיין מוביל בתחום שבו ההבדל נמדד בבירור, עם קנס יחס אות־רעש של \textenglish{0.26--2.53 dB}, התלוי ביעד השגיאה. אימות הממסר הנלמד בתקציב קבוע מגיע עד \textenglish{16 dB}; הזנב שמעבר לכך לא אומת. בסריקת פיילוטים נפרדת בערוץ מורכב ב־\textenglish{10 dB}, הממסר הנלמד מפגר אחרי הגילוי נעזר־הפיילוטים בעשרים פיילוטים ומוביל בעשרה. ממצאים אלה נוגעים למשפחות הפגמים שנכללו באימון, ולא למשפחות שלא נראו בו.

מחקר זיכרון נפרד בקישור \textenglish{QPSK} מקודד משתמש בממסר בן 220 פרמטרים וב־208 פעולות כפל־צבירה לסימבול. עלותו החישובית חוצה את הספירה המוצהרת של \textenglish{MLSE} בעל מלוא המצבים בין שלושה לחמישה מקדמים, בעוד הגלאי הקלאסי שומר על שיעור שגיאה נמוך יותר. מחקרי הממסור המקודד מראים הבדלים התלויים בתצורה ובמדדי הפענוח שמומשו, ולא מנגנון כללי ההופך פענוח בממסר לנחות. בקרת \textenglish{BCJR} מודדת שגיאות ביציאת הממסר בלבד.

הראיות תומכות ביתרונות מוגבלים ביחס לשרשראות הקלאסיות המסוימות שנבדקו. הן אינן מוכיחות עליונות אוניברסלית של ממסר נלמד, הכללה מחוץ למשפחות האימון או חיסכון בהשהיה ובאנרגיה בחומרה.
\end{justify}
\end{otherlanguage}
